\renewcommand{\Anno}{1997}

\renewcommand{\NumeroProblema}{6}

\renewcommand{\Testo}{
La figura sottostante rappresenta la piantina di un quartiere,
costituito da un reticolo ortogonale di strade equispaziate.
	\begin{center}
		\psset{unit=7mm}
		\pspicture(0,0)(3,3)
	
		\psset{linecolor=blue}
		\psline{-}(1,0)(1,3)
		\psline{-}(2,0)(2,3)
		\psline{-}(0,1)(3,1)
		\psline{-}(0,2)(3,2)
		\psline{-}(0,0)(3,0)(3,3)(0,3)(0,0)
	
		\psset{linecolor=black}
		\psdots[dotscale=1.5,dotstyle=square](3,3)
		\psdots[dotscale=1.5,dotstyle=o](0,0)
	
		\endpspicture
	\end{center}

Tizio abita al vertice Nord-Est (indicato con un quadratino), ed ha
l'ufficio al vertice Sud-Ovest (indicato con un cerchietto).  Caio
abita dove Tizio ha l'ufficio ed ha l'ufficio dove abita Tizio.

Tizio e Caio escono ogni mattina alla stessa ora e si recano al lavoro
camminando alla stessa velocit\`{a}: Tizio, ad ogni incrocio, procede
o verso Sud o verso Ovest; Caio, ad ogni incrocio, procede o verso
Nord o verso Est.  Ad ogni incrocio la scelta fra le possibili strade
che verificano la richiesta precedente \`{e} equiprobabile.

Determinare la probabilit\`{a} che Tizio e Caio si incontrino.
}

\renewcommand{\Risposta}{$41/100$}

\renewcommand{\Aiutino}{
In quali punti pu\`{o} avvenire l'incontro? Quanti sono in totale i 
possibili cammini? Quanti sono quindi i casi possibili? Quanti sono i 
cammini che passano per i possibili punti d'incontro? Quanti sono 
quindi i casi favorevoli?
}

\renewcommand{\SoluzioneUfficiale}{
Assumendo convenzionalmente che ogni quadratino della griglia abbia 
lato unitario, tutti i possibili percorsi di Tizio e Caio hanno 
lunghezza 6. Procedendo alla stessa velocit\`{a}, l'eventuale incontro 
potr\`{a} avvenire solo a distanza 3 dalle rispettive abitazioni, 
dunque solo nei 4 punti indicati nella figura con un triangolino.
\begin{center}
	\psset{unit=7mm}
	\pspicture(0,0)(3,3)

	\psset{linecolor=blue}
	\psline{-}(1,0)(1,3)
	\psline{-}(2,0)(2,3)
	\psline{-}(0,1)(3,1)
	\psline{-}(0,2)(3,2)
	\psline{-}(0,0)(3,0)(3,3)(0,3)(0,0)

	\psset{linecolor=black}
	\psdots[dotscale=1.5,dotstyle=square](3,3)
	\psdots[dotscale=1.5,dotstyle=o](0,0)
	\psdots[dotscale=1.5,dotstyle=triangle](0,3)(1,2)(2,1)(3,0)

	\endpspicture
\end{center}

Determiniamo quanti sono i cammini possibili per Tizio e per Caio.
Ogni cammino consiste in 6 tratti, di cui 3 orizzontali e 3 verticali.
Non \`{e} difficile quindi notare che i possibili cammini sono in
corrispondenza biunivoca con le sequenze di 6 lettere, di cui 3 sono V
e 3 sono O, quindi con i modi di scegliere 3 oggetti su 6, e sono
pertanto
$$\binom{6}{3}=20.$$

Determiniamo ora quanti sono i possibili cammini che passano per
ciascuno dei triangolini.  Per quelli ``esterni'' (ai vertici
Nord-Ovest e Sud-Est del quadrato) vi \`{e} un unico cammino.
Consideriamo dunque i 2 triangolini ``interni''.  \`{E} facile vedere
che, partendo ad esempio dal vertice Sud-Ovest ci sono 3 modi di
arrivarci, poi altri 3 modi per andare al vertice Nord-Est.  In tutto
ci sono quindi $3\cdot 3=9$ cammini che passano per ciascuno dei
triangolini interni.

Calcoliamo ora la probabilit\`{a} richiesta come rapporto tra casi 
favorevoli e casi possibili.

I casi possibili sono tutte le coppie (cammino di Tizio, cammino di 
Caio) e sono pertanto $20^{2}$. I casi favorevoli sono, tra le coppie 
precedenti, tutte e sole quelle in cui entrambi i cammini passano per 
uno stesso triangolino, e cio\`{e} $1^{2}$ per quelli esterni e 
$9^{2}$ per quelli interni.

Pertanto la probabilit\`{a} richiesta \`{e}
$$\frac{1^{2}+9^{2}+9^{2}+1^{2}}{20^{2}}=\frac{162}{400}=\frac{41}{100},$$
e cio\`{e} del 41\%.
}

\renewcommand{\SoluzioniAlternative}{
}

\renewcommand{\Osservazioni}{\TitoloOss{1} \`{E} possibile contare i
cammini senza ricorrere all'interpretazione combinatoria ed ai
coefficienti binomiali.  Supponiamo di partire dal vertice Sud-Ovest e
segniamo in ogni incrocio quanti sono i possibili cammini che vi
passano.

\`{E} chiaro che tale numero \`{e} 1 per tutti gli incroci sul lato
Sud e sul lato Ovest.  Per ogni altro incrocio, il numero dei cammini
che vi passano \`{e} uguale alla somma dei cammini passanti per
l'incrocio immediatamente ad Ovest pi\`{u} quelli passanti per
l'incrocio immediatamente a Sud.  Mediante questa regola possiamo
quindi pian piano riempire le rimanenti caselle, fino al vertice
Nord-Est.

\begin{center}
	\psset{unit=3mm}
	\pspicture(-1,-1)(13,13)
	
	\multips(0,0)(0,4){4}{\multips(0,0)(4,0){3}{
	\psline[linecolor=blue]{->}(1,0)(3,0)}}
	
	\multips(0,0)(0,4){3}{\multips(0,0)(4,0){4}{
	\psline[linecolor=blue]{->}(0,1)(0,3)}}
	
	\multips(0,0)(0,4){4}{\multips(0,0)(4,0){4}{\pscircle(0,0){1}}}
	
	\rput(0,0){1}
	\rput(4,0){1}
	\rput(8,0){1}
	\rput(12,0){1}
	
	\rput(0,4){1}
	\rput(4,4){2}
	\rput(8,4){3}
	\rput(12,4){4}
	
	\rput(0,8){1}
	\rput(4,8){3}
	\rput(8,8){6}
	\rput(12,8){10}
	
	\rput(0,12){1}
	\rput(4,12){4}
	\rput(8,12){10}
	\rput(12,12){20}
	
	\endpspicture
\end{center}


Si verifica cos\`{\i} che i possibili cammini da Sud-Ovest a Nord-Est sono 20, 
cos\`{\i} come i cammini da Sud-Ovest fino ad un triangolino interno sono 3.

\TitoloOss{2} Si noti l'analogia tra il procedimento descritto nella
precedente osservazione e la costruzione del triangolo di Tartaglia.
}

\renewcommand{\Commenti}{
}